% 数学ノート集・共通テンプレート v2（積分論2026年度第1部の版面を参考）
% A5判。章内で定義・定理・例・注意等の番号を共通化。
% コンパイル： platex -interaction=nonstopmode math_note_template_v2.tex (2回)
%              dvipdfmx math_note_template_v2.dvi
\documentclass[a5paper,10pt,onecolumn]{jsarticle}
\usepackage[dvipdfmx,top=22mm,bottom=20mm,inner=19mm,outer=18mm,headheight=13pt,headsep=6mm]{geometry}
\usepackage{amsmath,amssymb,amsthm,mathtools}
\usepackage{enumitem,fancyhdr,titlesec}
\usepackage[dvipdfmx,bookmarks=true,colorlinks=false]{hyperref}
\usepackage{pxjahyper}

% ----- 差し替える文書情報 -----
\newcommand{\NoteTitle}{Stone--Weierstrassの定理}
\newcommand{\NoteSeries}{数学ノート集}

% ----- 文字・組版 -----
% 参考PDFの質素な明朝＋Computer Modernの印象を引き継ぎながら，
% A5判で窮屈にならないよう本文の行間を少し確保する．
\linespread{1.13}
\setlength{\parindent}{1zw}
\setlength{\parskip}{0pt}
\setlength{\abovedisplayskip}{6pt plus 2pt minus 2pt}
\setlength{\belowdisplayskip}{6pt plus 2pt minus 2pt}
\setlength{\abovedisplayshortskip}{3pt plus 2pt minus 1pt}
\setlength{\belowdisplayshortskip}{4pt plus 2pt minus 1pt}
\setlist[enumerate]{topsep=3pt,itemsep=1pt,parsep=0pt,leftmargin=2.4zw}
\titleformat{\section}{\normalfont\normalsize\bfseries}{\thesection}{0.8zw}{}
\titlespacing*{\section}{0pt}{2.0ex plus .4ex}{.9ex}
\titleformat{\subsection}{\normalfont\normalsize\bfseries}{\thesubsection}{0.7zw}{}
\titlespacing*{\subsection}{0pt}{1.5ex plus .3ex}{.6ex}
\raggedbottom

% ----- 柱とページ番号：ページ上部、罫線なし -----
\fancypagestyle{mathnote}{%
  \fancyhf{}%
  \fancyhead[L]{\footnotesize \NoteSeries\ （\NoteTitle）}%
  \fancyhead[R]{\footnotesize \thepage}%
  \renewcommand{\headrulewidth}{0pt}%
  \renewcommand{\footrulewidth}{0pt}%
}
\pagestyle{mathnote}

% ----- 節内通し番号の定義・定理群 -----
\newtheoremstyle{mathnotetheorem}
 {4.5pt}{4.5pt}{\normalfont}{}%
 {\bfseries}{}{0.65zw}%
 {\thmname{#1}\thmnumber{ #2}\thmnote{（#3）}}
\theoremstyle{mathnotetheorem}
\newtheorem{theorem}{定理}[section]
\newtheorem{definition}[theorem]{定義}
\newtheorem{lemma}[theorem]{補題}
\newtheorem{proposition}[theorem]{命題}
\newtheorem{corollary}[theorem]{系}
\newtheorem{example}[theorem]{例}
\newtheorem{remark}[theorem]{注意}
\renewcommand{\proofname}{証明}
\renewcommand{\qedsymbol}{$\square$}
\makeatletter
\renewenvironment{proof}[1][\proofname]{%
  \par\pushQED{\qed}\normalfont\topsep4\p@\@plus2\p@\relax
  \trivlist\item[\hskip\labelsep\bfseries #1]\ignorespaces
}{\popQED\endtrivlist\@endpefalse}
\makeatother

% ----- 使うノートで適宜拡張する記号 -----
\newcommand{\R}{\mathbb{R}}
\newcommand{\norm}[1]{\lVert #1\rVert}
\newcommand{\abs}[1]{\lvert #1\rvert}

\begin{document}
\section*{読了後にできること}
本稿を読むための前提は，一様収束，連続関数，Stone--Weierstrassの定理．
本稿の具体的な到達課題は，具体的な連続関数に近づく多項式列を構成し，誤差を計算する．

\section{連続関数の代数と一様近似}

$K$ を空でないコンパクトHausdorff空間とする．
$C(K,\mathbb F)$ は $K$ 上の $\mathbb F$ 値連続関数全体であり，
一様ノルム $\norm f_\infty=\max_{x\in K}|f(x)|$ を入れる．
ここで $\mathbb F=\mathbb R$ または $\mathbb C$ とする．

\begin{definition}
$\mathcal A\subset C(K,\mathbb F)$ が\emph{部分代数}であるとは，
$\mathcal A$ が $\mathbb F$ 線形部分空間であり，
$f,g\in\mathcal A$ なら $fg\in\mathcal A$ であることをいう．
\emph{定数を含む}とは定数関数 $1$ が $\mathcal A$ に属することをいう．
$\mathcal A$ が\emph{点を分離する}とは，
$x\ne y$ に対して $f(x)\ne f(y)$ となる
$f\in\mathcal A$ が存在することをいう．
\end{definition}

\begin{lemma}\label{lem:closed-algebra}
$\mathcal A\subset C(K,\mathbb F)$ が部分代数なら，
その一様閉包 $\overline{\mathcal A}$ も部分代数である．
定数を含む性質と点を分離する性質は閉包に移っても保たれる．
\end{lemma}
\begin{proof}
$f_n\to f$，$g_n\to g$ が一様収束するとき，
$f_n+g_n\to f+g$ と $af_n\to af$ は一様収束する．
また $(f_n)$ は一様有界なので
\[
 \norm{f_ng_n-fg}_\infty
 \le\norm{f_n}_\infty\norm{g_n-g}_\infty
     +\norm{f_n-f}_\infty\norm g_\infty\to0.
\]
従って閉包は部分代数となる．
定数関数が属することは明らかであり，
点を分離する元はもともと $\mathcal A$ に属するから
閉包においても存在する．
\end{proof}

\section{Bernstein多項式による近似}

\begin{definition}
$f\in C([0,1],\mathbb R)$ と $n\ge1$ に対し，
Bernstein多項式を
\[
 B_nf(x)=\sum_{k=0}^nf(k/n){n\choose k}x^k(1-x)^{n-k}
\]
と定義する．
\end{definition}

\begin{lemma}[二項係数による恒等式]\label{lem:binomial}
$0\le x\le1$ とし，
$w_{n,k}(x)={n\choose k}x^k(1-x)^{n-k}$ と置く．
このとき
\begin{align*}
 \sum_{k=0}^n w_{n,k}(x)&=1,\\
 \sum_{k=0}^n\frac kn w_{n,k}(x)&=x,\\
 \sum_{k=0}^n\left(\frac kn-x\right)^2w_{n,k}(x)
 &=\frac{x(1-x)}n\le\frac1{4n}.
\end{align*}
\end{lemma}
\begin{proof}
最初の等式は二項定理による．
$k{n\choose k}=n{n-1\choose k-1}$ を用いれば
\[
 \sum_{k=0}^nk\,w_{n,k}(x)
 =nx\sum_{j=0}^{n-1}{n-1\choose j}x^j(1-x)^{n-1-j}
 =nx
\]
である．同様に
$k(k-1){n\choose k}=n(n-1){n-2\choose k-2}$ から
$\sum_k k(k-1)w_{n,k}=n(n-1)x^2$ を得る．
$k^2=k(k-1)+k$ を使って二次の式を展開すれば
$\sum_k(k-nx)^2w_{n,k}=nx(1-x)$ となる．
$n=1$ のときも直接計算で同じ式が成立する．
\end{proof}

\begin{theorem}[Weierstrassの近似定理]\label{thm:weierstrass}
任意の $f\in C([0,1],\mathbb R)$ に対して
$B_nf\to f$ は $[0,1]$ 上で一様収束する．
従って任意の有界閉区間上の連続実数値関数は，
実係数多項式によって一様近似できる．
\end{theorem}
\begin{proof}
$f$ は一様連続だから，
任意の $\varepsilon>0$ に対して
$|u-v|<\delta$ なら
$|f(u)-f(v)|<\varepsilon/2$ となる $\delta>0$ を取る．
$M=\norm f_\infty$ と置けば，
補題\ref{lem:binomial}から
\begin{align*}
 |B_nf(x)-f(x)|
 &\le\sum_{|k/n-x|<\delta}
  |f(k/n)-f(x)|w_{n,k}(x)\\
 &\quad+2M\sum_{|k/n-x|\ge\delta}w_{n,k}(x)\\
 &\le\frac{\varepsilon}{2}
 +\frac{2M}{\delta^2}
   \sum_{k=0}^n(k/n-x)^2w_{n,k}(x)\\
 &\le\frac{\varepsilon}{2}+\frac{M}{2n\delta^2}.
\end{align*}
最後の項は $n$ を十分大きくすれば
$\varepsilon/2$ 未満となる．
この評価は $x\in[0,1]$ に依存しないため，一様収束を得る．
一般の閉区間 $[a,b]$ については
$x=a+(b-a)t$ と変数変換すればよい．
$a=b$ は定数多項式で自明である．
\end{proof}

\begin{corollary}\label{cor:absolute}
任意の $M>0$ と $\varepsilon>0$ に対し，
実係数多項式 $p$ で
\[
 \sup_{|t|\le M}\bigl|p(t)-|t|\bigr|<\varepsilon
\]
を満たすものが存在する．
\end{corollary}
\begin{proof}
$[-M,M]$ 上の連続関数 $t\mapsto|t|$ に
定理\ref{thm:weierstrass}を適用する．
\end{proof}

\section{格子作用素の構成}

\begin{lemma}\label{lem:lattice}
$\mathcal A\subset C(K,\mathbb R)$ を
定数を含む部分代数とし，
$B=\overline{\mathcal A}$ とする．
このとき任意の $f,g\in B$ に対して
$|f|$，$\max\{f,g\}$，$\min\{f,g\}$ は $B$ に属する．
\end{lemma}
\begin{proof}
$B$ は補題\ref{lem:closed-algebra}により閉部分代数である．
$M=\norm f_\infty+1$ と置く．
系\ref{cor:absolute}により，
$[-M,M]$ 上で $|t|$ に一様収束する実係数多項式列
$(p_n)$ が存在する．
$B$ は定数を含む代数なので $p_n(f)\in B$ である．
\[
 \norm{p_n(f)-|f|}_\infty
 \le\sup_{|t|\le M}|p_n(t)-|t||\to0
\]
だから $|f|\in B$ である．
次に
\[
 \max\{f,g\}=\frac{f+g+|f-g|}{2},\qquad
 \min\{f,g\}=\frac{f+g-|f-g|}{2}
\]
より両者は $B$ に属する．
\end{proof}

\begin{lemma}[二点補間]\label{lem:interpolation}
$\mathcal A\subset C(K,\mathbb R)$ が
定数を含み点を分離するとする．
任意の $h\in C(K,\mathbb R)$ と $x,y\in K$ に対し，
$g\in\mathcal A$ で
$g(x)=h(x)$，$g(y)=h(y)$ を満たすものが存在する．
\end{lemma}
\begin{proof}
$x=y$ なら $g$ を定数関数 $h(x)$ とすればよい．
$x\ne y$ なら点分離性により
$a\in\mathcal A$ で $a(x)\ne a(y)$ を取れる．
\[
 g(z)=h(x)+\frac{h(y)-h(x)}{a(y)-a(x)}
        (a(z)-a(x))
\]
とすれば $g\in\mathcal A$ であり，
必要な二点での等式を満たす．
\end{proof}

\section{実Stone--Weierstrassの定理}

\begin{theorem}[実Stone--Weierstrass]\label{thm:real-sw}
$K$ をコンパクトHausdorff空間とする．
$\mathcal A\subset C(K,\mathbb R)$ が
定数関数を含む部分代数であって点を分離するなら，
$\mathcal A$ は $C(K,\mathbb R)$ に一様稠密である．
\end{theorem}
\begin{proof}
$B=\overline{\mathcal A}$ と置く．
補題\ref{lem:lattice}より $B$ は最大・最小について閉じている．
任意の $h\in C(K,\mathbb R)$ と $\varepsilon>0$ を固定する．

まず $x\in K$ を固定する．
各 $y\in K$ に対して補題\ref{lem:interpolation}から
$g_{x,y}\in\mathcal A$ で
$g_{x,y}(x)=h(x)$，
$g_{x,y}(y)=h(y)$ となるものを取る．
連続性により
\[
 U_{x,y}=\{z\in K:g_{x,y}(z)<h(z)+\varepsilon\}
\]
は $y$ を含む開集合である．
$(U_{x,y})_{y\in K}$ は $K$ の開被覆なので，
ある $y_1,\dots,y_m$ について
$K=\bigcup_{j=1}^mU_{x,y_j}$ となる．
\[
 g_x=\min_{1\le j\le m}g_{x,y_j}\in B
\]
と置くと，任意の $z\in K$ について
$g_x(z)<h(z)+\varepsilon$ であり，
$g_x(x)=h(x)$ である．
従って
$V_x=\{z\in K:g_x(z)>h(z)-\varepsilon\}$
は $x$ を含む開集合である．

$(V_x)_{x\in K}$ は $K$ を被覆するから，
ある $x_1,\dots,x_N$ について
$K=\bigcup_{i=1}^NV_{x_i}$ とできる．
\[
 g=\max_{1\le i\le N}g_{x_i}\in B
\]
と置けば，
任意の $z\in K$ で少なくとも一つの $i$ が
$z\in V_{x_i}$ を満たすため
$g(z)>h(z)-\varepsilon$ となる．
一方，全ての $i$ で
$g_{x_i}(z)<h(z)+\varepsilon$ であるから
$g(z)<h(z)+\varepsilon$ となる．
従って $\norm{g-h}_\infty\le\varepsilon$ である．
$\varepsilon>0$ は任意なので $h\in B$ となり，
$B=C(K,\mathbb R)$ を得る．
\end{proof}

\section{複素Stone--Weierstrassの定理}

\begin{definition}
部分代数 $\mathcal A\subset C(K,\mathbb C)$ が
\emph{自己共役}であるとは，
$f\in\mathcal A$ に対して
$\overline f\in\mathcal A$ が成り立つことをいう．
ここで $\overline f(x)=\overline{f(x)}$ とする．
\end{definition}

\begin{theorem}[複素Stone--Weierstrass]\label{thm:complex-sw}
$K$ をコンパクトHausdorff空間とする．
$\mathcal A\subset C(K,\mathbb C)$ が，
定数を含み，点を分離し，自己共役な部分代数なら，
$\mathcal A$ は $C(K,\mathbb C)$ に一様稠密である．
\end{theorem}
\begin{proof}
$\mathcal A_\mathbb R=\{f\in\mathcal A:
 f(K)\subset\mathbb R\}$ と置く．
これは定数を含む $C(K,\mathbb R)$ の実部分代数である．
$x\ne y$ のとき，$\mathcal A$ の点分離性により
$g(x)\ne g(y)$ となる $g\in\mathcal A$ が存在する．
実部または虚部の少なくとも一方は $x,y$ で異なる．
自己共役性により
\[
 \operatorname{Re}g=\frac{g+\overline g}{2},\qquad
 \operatorname{Im}g=\frac{g-\overline g}{2i}
\]
はいずれも $\mathcal A_\mathbb R$ に属する．
従って $\mathcal A_\mathbb R$ は点を分離する．
定理\ref{thm:real-sw}より
$\mathcal A_\mathbb R$ は $C(K,\mathbb R)$ に稠密である．

任意の $h\in C(K,\mathbb C)$ に対し，
$\operatorname{Re}h$ と $\operatorname{Im}h$ をそれぞれ
$\mathcal A_\mathbb R$ の元 $u_n,v_n$ で一様近似すれば，
$u_n+iv_n\in\mathcal A$ かつ
$\norm{u_n+iv_n-h}_\infty\to0$ を得る．
\end{proof}

\begin{example}[自己共役条件が必要な例]
$K=\{z\in\mathbb C:|z|=1\}$ とし，
$\mathcal A$ を複素多項式 $p(z)$ の $K$ への制限全体とする．
$\mathcal A$ は定数を含み点を分離する部分代数であるが，
$C(K,\mathbb C)$ に一様稠密でない．
\end{example}
\begin{proof}
$h(z)=\overline z=1/z$（$z\in K$）を考える．
$\mathcal A$ の元 $p$ には原始関数があるから
\[
 \int_{|z|=1}p(z)\,dz=0
\]
である．しかし，正向きの単位円周について
\[
 \int_{|z|=1}\overline z\,dz
 =\int_{|z|=1}\frac{dz}{z}=2\pi i.
\]
もし $p_n\to h$ が $K$ 上で一様収束するなら，
曲線の長さによる積分評価により
$\int p_n\,dz\to\int h\,dz$ となって矛盾する．
\end{proof}

\section{具体的な近似定理への応用}

\begin{corollary}[多変数多項式の稠密性]
$K\subset\mathbb R^d$ を空でないコンパクト集合とする．
$\mathbb R^d$ 上の実係数多項式の $K$ への制限は
$C(K,\mathbb R)$ に一様稠密である．
\end{corollary}
\begin{proof}
多項式の制限全体は定数を含む部分代数である．
異なる $x,y\in K$ に対してある座標 $j$ が
$x_j\ne y_j$ を満たすため，
座標関数 $z\mapsto z_j$ により点を分離する．
定理\ref{thm:real-sw}を適用する．
\end{proof}

\begin{corollary}[三角多項式の稠密性]
連続な $2\pi$ 周期複素数値関数 $f$ に対し，
三角多項式
\[
 P(t)=\sum_{k=-N}^Nc_ke^{ikt}
\]
で $f$ を一様近似できる．
\end{corollary}
\begin{proof}
円周 $K=\{z:|z|=1\}$ 上の関数
$z\mapsto z^k$（$k\in\mathbb Z$）の有限線形結合は
定数を含む部分代数であり，
$z\mapsto z$ が点を分離する．
$\overline{z^k}=z^{-k}$ が成り立つので自己共役である．
定理\ref{thm:complex-sw}を適用し，
$z=e^{it}$ と置けば結論を得る．
\end{proof}

\begin{remark}
Stone--Weierstrassの定理は，近似する函数の具体的な係数を
与える定理ではない．一方，Bernstein多項式は
一変数の場合に明示的な近似列を与える．
両者は一様近似の存在と構成という異なる役割を持つ．
\end{remark}

\section{到達課題の完全解答}
\noindent\textbf{Bernstein多項式で一様近似する．}
\begin{proposition}
$[0,1]$ 上の $f(x)=x^2$ に対するBernstein多項式
\[
B_nf(x)=\sum_{k=0}^n
\left(\frac kn\right)^2\binom nkx^k(1-x)^{n-k}
\]
は
\[
B_nf(x)=x^2+\frac{x(1-x)}n
\]
である．従って
$\|B_nf-f\|_\infty=1/(4n)$ である．
\end{proposition}
\begin{proof}
$K\sim\operatorname{Bin}(n,x)$ とすると
$B_nf(x)=E[(K/n)^2]$ である．
二項分布の平均と分散を使うと
\[
E[(K/n)^2]
=\frac{\operatorname{Var}(K)+(EK)^2}{n^2}
=x^2+\frac{x(1-x)}n.
\]
$[0,1]$ 上の $x(1-x)$ の最大値は
$x=1/2$ での $1/4$ なので誤差が従う．
\end{proof}
\noindent\textbf{到達確認．}
$f(x)=x$ に対する $B_nf$ を求めよ．
\par\noindent\textbf{解答．}
$B_nf(x)=E[K/n]=x$ であり，
一次関数はどの次数でも正確に再現される．
\end{document}
